%%%%%
%
\section{Rotation vers le rep�re port� par le vent moyen}

\itemstyle{Op\'erateur~:} {\sf Rotat\_rep\_vent\_moy}
\index{Rotat\_rep\_vent\_moy}

%%%%%OBJET%%%%
\itemstyle{Objet~:}
Calcul des 3 composantes du vecteur-vent dans 
un rep�re port� par le vent moyen, qui v�rifie : 
\[
\begin{array}{c}
\overline{u} = U\\
\overline{v} = 0\\
\overline{w} = 0\\
\overline{v'w'} = 0\\
\end{array}
\]

  En pratique, on minimise (mais n'annule pas toujours) le 
terme $\overline{v'w'}$.

%%%%%ARGUMENTS%%%%
\bigskip
%%%%%ENTREES%%%%
\begin{tabarg}
\multicolumn{3}{l}{\titretab {Entr\'ees}}\\
\hline
$u$ & Vecteur, {\tt float} & premi�re composante horizontale du vent (any)\\
$v$ & Vecteur, {\tt float} & seconde composante horizontale du vent (m�me unit� que $u$)\\
$w$ & Vecteur, {\tt float} & composante verticale du vent (m�me unit� que $u$)\\
\end{tabarg}

\medskip
%%%%%SORTIES%%%%
\begin{tabarg}
\multicolumn{3}{l}{\titretab {Sorties}}\\
\hline
$u_r$ & Vecteur, {\tt float} & composante du vent sur l'axe du vent moyen (m�me unit� que $u$)\\
$v_r$ & Vecteur, {\tt float} & composante du vent sur le premier axe perpendiculaire au vent moyen (m�me unit� que $u$)\\
$w_r$ & Vecteur, {\tt float} & composante du vent sur le second axe perpendiculaire au vent moyen (m�me unit� que $u$)\\
\end{tabarg}


%%%%%FORMULE%%%%
\itemstyle{Formule}

%%Rotat_rep_vent_moy
Application d'une rotation destin�e � annuler
les composantes transverses moyennes du vent :
\[
  \left(
\begin{array}{c}
u_r \\ v_r \\ w_r \\
\end{array}
  \right)
=
  \left(
\begin{array}{ccc}
\cos\epsilon \cos\theta & \sin\epsilon \cos\theta & \sin\theta \\
-\sin\epsilon & \cos\epsilon & 0 \\
-\cos\epsilon \sin\theta & -\sin\epsilon \sin\theta & \cos\theta \\
\end{array}
  \right)
.
  \left(
\begin{array}{c}
u \\ v \\ w \\
\end{array}
  \right)
\]

avec $\cos\epsilon=\frac{\overline{u}}{\sqrt{\overline{u}^2+\overline{v}^2}}$,
    $\sin\epsilon=\frac{\overline{v}}{\sqrt{\overline{u}^2+\overline{v}^2}}$,
    $\cos\theta=\frac{\sqrt{\overline{u}^2+\overline{v}^2}}
                     {\sqrt{\overline{u}^2+\overline{v}^2+\overline{w}^2}}$,
et  $\sin\theta=\frac{\overline{w}}
                     {\sqrt{\overline{u}^2+\overline{v}^2}+\overline{w}^2}$ 
(avec $\theta$ limit� � 20 degr�s).

Puis d�termination d'un angle $\beta$. minimisant $\overline{v'w'}$,
par r�solution it�rative de $2K\sin^2\beta+...= 0$



%%%%%COMMENTAIRE%%%%
\itemstyle{ Commentaire~: }

Le rep�re ($u$, $v$, $w$) d'origine peut �tre g�ographique ou li� � l'instrument. 
Il doit �tre orthonorm�.
%%%% Question en suspens : doit-il �tre direct ?


%%%%%AUTEUR%%%%
\itemstyle{Auteur~:}
CNRM/GMEI/TRAMM

%%%%%REFERENCE%%%%
\itemstyle{R\'ef\'erences~:}
Pages 26 � 28 de la note de P. Bessemoulin sur la
d�termination des flux de surface \cite{GMEI-6}.
